%% kapa: exam=1 round=26 year=2015 period=1 form=A subject=civil_law subject_ko=민법 seq=5 number=5
%% source: Q-Net 감정평가사 기출문제 (cst00309) artlSeq=5138001 file=extracted/1008/5138001/제26회_감정평가사 1차_기출/제26회_감정평가사 1차_기출/1교시_감평_A형.pdf
%% conversion: kordoc 4.24.51 -> tools/gen1.py (auto); review=proofread-sonnet; flags=none
\begin{question}{5}
대리에 관한 설명으로 옳은 것은? (다툼이 있으면 판례에 따름)
\begin{choices}
  \item 임의대리인은 행위능력자여야 한다.
  \item 대리인의 법률행위의 효과는 본인에게 귀속되므로 의사표시의 하자의 유무는 본인을 기준으로 판단한다.
  \item 복대리인은 대리인이 선임한 자로서 본인의 대리인이 아니다.
  \item 대리인은 본인의 허락을 얻어 본인을 위하여 자기와 법률행위를 할 수 있다.
  \item 권한을 넘은 표현대리의 규정은 법정대리에는 적용되지 않는다.
\end{choices}
\end{question}
\answer{4}
